\CatchFileDef{\singlelenashr}{single-res-table-shr-Lena.tex}{}

\section{Proposed Method}
\label{sec:ours}

In this section, we propose a new reversible data hiding in encrypted images method base on~\cite{Yin2020} with three improvements that have a significant increment on embedding rate (ER) while still complying with the typical requirements described in Section~\ref{sec:intro}. 



\subsection{Apply the Huffman tree building algorithm}
\label{sec:improvetree}

We observed that Yin et al.\ introduced their method based on a predefined Huffman tree 
without clarifying the pros and cons behind this design decision. 
We found that this fixed predefined tree topology constrains the true ability of Huffman coding. 
When we apply the standard Huffman coding algorithm
, its full power can be released~\cite{IRE1952, Leeuwen2008a}.

The tree topology of 
\emph{Lena} $512 \times 512$ is shown in Figure~\ref{fig:lenatree} and the statistic data is given in Table~\ref{tab:shrlena}. 
Since in our proposal the tree topology is no longer a preshared knowledge between the \emph{Content-Owner}, the \emph{Data-hider} and the \emph{Receiver}, we have to extend $S^{aux}$ to convey this information. 
 
A Huffman tree of 9 leaves has $2 \times 9 -1 = 17$ nodes, each needs one bit to represent its role, where 1 is defined for an internal node and 0 for a leaf~\cite{Pu2006}. 
Each leaf has a value $0 \leq v \leq 8$, thus four bits are enough to store each leaf value, i.e., $4 \times 9 = 36$ bits. 
Hence we occupy 53 bits in $S^{aux}$ to store the tree topology and leaf values.

\begin{figure}
	\scriptsize
	\Tree[ 
	[.0 [.0 \{00\}\\5 ]
	[.1 [.0 \{010\}\\7 ] [.1 \{011\}\\8 ]]]
	[.1 [.0 [.0 \{100\}\\3 ]
	[.1 [.0 \{1010\}\\2 ]
	[.1 [.0 \{10110\}\\1 ] [.1 \{10111\}\\0 ]]]]
	[.1 [.0 \{110\}\\6 ] [.1 \{111\}\\4 ]]]] 
	\caption{The Huffman coding tree of \emph{Lena} $512 \times 512$.}
	\label{fig:lenatree}
\end{figure}


\begin{table}
	%	\scriptsize
	\centering
	\caption{\emph{Lena} $512 \times 512$ using a different Huffman tree.}
	
	\begin{tabular}{crlcccr}
		\toprule 
		$v$ &   \multicolumn{1}{c}{$f_v$} &     \multicolumn{1}{c}{$encode$}  &  $cap$ &    $\vert encode \vert$    & $payload$  & $\Delta$ \\
		\midrule  
		\singlelenashr
		\bottomrule   
	\end{tabular}
	
	\label{tab:shrlena}
\end{table}

\subsection{Reduce Reference Bytes}
\label{sec:improveref}

In the method of Yin et al.,
the bytes in the first row and the first column, i.e.\ $REF_e$, must be losslessly preserved, and this is the reason why they must concatenate $REF_e$ after $T^\prime$ as part of $S^{aux}$ and embed them before embedding any secret information 
. 

Though this design does not provide any negative effect since $REF_e$ are $neutral$ base on the fact that every byte in $REF_e$ provides 8-bit capacity while also occupying 8 bits to embed itself, with one more byte being defined as $REF_e$, there is one less byte that can be used for information hiding. 

As shown in~\eqref{eq:pe2}, 
for the bytes in the first row, we define their predicted error $pv$ as exclusive-or with the previous byte. 
Similar logic applies to those in the first column. 
\begin{equation}
	\label{eq:pe2}
	pe(i, j) = 
	\begin{cases}
		\text{undefined},             & \text{if } i = 0 \wedge  0 \leq j < |REF_e| \\
		x_{i, j - 1} \oplus x_{i, j}, & \text{if } i = 0 \\
		x_{i - 1, j} \oplus x_{i, j}, & \text{if } j = 0 \\
		pv(i, j) \oplus x_{i, j}, 	  & \text{if } 1 \leq i < h \wedge 1 \leq j < w\,. \\
	\end{cases} 
\end{equation}

Theoretically, the top-left corner pixel $x_{0, 0}$ is the only critical byte that must be losslessly preserved, but for the following reasons, we define $REF_e$ more than this one byte.
Aside from the reason described in Section~\ref{sec:improvetree} that we need 53 bits to store the Huffman tree, 
more importantly, we need some more bytes in $REF_e$ to avoid violating a loop invariant during 
their information extraction process~\cite{Furia2014}.

There are two indexes in 
their information extraction process, 
i.e., the $(d_x, d_y)$ 
and the $(i, j)$ for the loop.
For 
their information extraction process
to work properly, there is a never mentioned loop invariant  
that $(i, j) < (d_x, d_y)$ must be true, i.e., the index $(d_x, d_y)$ must be in a more advanced position in $I_w$ than $(i, j)$ for all time.

While $(i, j)$ is increased with every loop iteration, $(d_x, d_y)$ is only increased when the $buf$ 
contains a proper Huffman code, which is affected by the capacity of each proceeded $t_{i, j}$.
If the distribution of $negative$ $t_{i, j}$ has a tendency to centralized in some part of $T$ and slows down the incremental speed of $(d_x, d_y)$, 
their information extraction process
will encounter exception and fail.

Increase $REF_e$ helps to avoid this problem. 
As can be seen 
in their information extraction process,
after extracting the information of Huffman tree, the rest of the $REF_e$ will be consumed to produce more $t_{(d_x, d_y)}$ and put $(d_x, d_y)$ at a more advanced position, thus reducing the chance of violating the loop invariant.
Though defining more $REF_e$ bytes helps to amend this issue, it also decreases the bytes that can be dedicated to information hiding.

Even though the minimum $REF_e$ requirement differs from image to image, for example, the Airplane $512 \times 512$ needs only 1 extra $REF_e$ beside the basic 53 bits used to convey the coding tree, while Baboon $512 \times 512$ asks for 16 more $REF_e$ bytes to safely go through 
their information extraction process, 
this loop invariant only exists in extracting process and remain unnoticed to the \emph{Content-Owner}.
We have to define a fixed and large enough $REF_e$ and preshare this knowledge with all three participators in the use case. 

According to our experiments, defining $REF_e$ equals to or larger than $7 + 63$ bytes is enough for all the test images we can acquire~\cite{UniversityofSouthernCalifornia2020, Bas2011, bows2007}, but we still retain the possibility that it might need to be extended for some more complex cover images.


















\subsection{Remove the redundant length column}
\label{sec:improvesize}


While we inspecting the method of Yin et al., we found the unreasonable design of constraining the size of $LEN$ to 20 bits, which is only suitable for images of size smaller than or equal to $512 \times 512$. 
When using an image of larger size, e.g., $1024 \times 1024$, 
it is understandable that the label map $T$ is of size $1023 \times 1023$ by definition, which is only slightly less than $2^{20}$, then the total $LEN$ of $T^\prime$ will sure be exceeded the size of 20 bits. 
In fact, according to the experimental data published by~\cite{Yin2020}, the image \emph{Man $1024 \times 1024$} has $LEN = 3,121,508$ bits, which requires at least 22 bits to represent it since $log_2(3,121,508) = 21.574$.
This can be amended by allocating enough space in the auxiliary information by~\eqref{eq:floor}.
\begin{equation}
	\label{eq:floor}
	\vert LEN \vert \; \leq \lfloor log_2((w - 1) \times (h - 1) \times C ) \rfloor + 1 , 
\end{equation}
where the constant $C$ is the Shannon entropy~\cite{Shannon1948} representing the lower bound of the number of bits required to encode each value $t \in T$ 
in the worst-case scenario when $F$ is a uniform distribution.

When we further analyze 
their information extraction process, 
we realized that the $LEN$ is redundant information. 
Since for both the secret extracting process
and the cover image restoring process,  
the label map $T$ is essential information, 
and the $LEN$ is a by-product of the process of re-constructing $T$ from $I_e^\prime$ or $I_w$, 
hence do not need to be explicitly conveyed in $S^{aux}$. 
In our experiments, it is safe to remove the column $LEN$ from $S^{aux}$ 
.
